% GENERATED FILE. Edit canonical lessons, then run npm run generate:books.
% canonical-source-hash: fnv1a64:dc9bfee83dc93dc6
% canonical-lessons: MR-C223-durust-karne, MR-C223-photo-kadhne, MR-C223-samjavne, MR-C223-bhasantar-karne, MR-C223-nirnay-ghene

\chapter{To repair, To take a photo, To explain, To translate, To make a decision}
\label{ch:mr-to-repair-to-take-a-photo-to-explain-to-translate-to-make-a-decision}
\hlchaptermodality{223}

\begin{chapteropening}
\textbf{By the end of this chapter:} \emph{I can say to repair, to take a photo, to explain, to translate and to make a decision.}

Complete the last lesson of chapter 223: I can say to repair, to take a photo, to explain, to translate and to make a decision.
\end{chapteropening}

\section[\texorpdfstring{\mr{दुरुस्त} \mr{करणे} (durust karṇe)}{दुरुस्त करणे (durust karṇe)}]{\mr{दुरुस्त} \mr{करणे} (durust karṇe) — to repair}

\label{lesson:MR-C223-durust-karne}



\begin{quote}
Before the new one: say the Marathi for a university, then the Marathi for a mistake.
\end{quote}

\subsection*{You'll want to know: \mr{दुरुस्त} \mr{करणे}}
\textbf{\mr{दुरुस्त} \mr{करणे}} — \emph{durust karṇe} — \textquotedblleft{}to repair\textquotedblright{}.

\begin{grammarlens}[title={What you've built}]
One more word for this chapter.
\end{grammarlens}

\subsection*{Your turn}
\begin{itemize}
\raggedright
  \item \emph{Say it:} \emph{durust karṇe}
  \item \emph{Say it:} \emph{durust karṇe}, once more
  \item \emph{Say it:} say \emph{durust karṇe}
  \item \emph{Recall:} say the Marathi for a university, then the Marathi for a mistake, then say \emph{durust karṇe} again
\end{itemize}

\begin{tcolorbox}[breakable,colback=teal!4,colframe=teal!35!black,title={Before you move on}]
What does \mr{दुरुस्त} \mr{करणे} mean? (\textquotedblleft{}To repair\textquotedblright{}.) Say it once more.
\end{tcolorbox}
\section[\texorpdfstring{\mr{फोटो} \mr{काढणे} (phoṭo kāḍhṇe)}{फोटो काढणे (phoṭo kāḍhṇe)}]{\mr{फोटो} \mr{काढणे} (phoṭo kāḍhṇe) — to take a photo}

\label{lesson:MR-C223-photo-kadhne}



\begin{quote}
Before the new one: say the Marathi for a mistake, then the Marathi for to repair.
\end{quote}

\subsection*{You'll want to know: \mr{फोटो} \mr{काढणे}}
\textbf{\mr{फोटो} \mr{काढणे}} — \emph{phoṭo kāḍhṇe} — \textquotedblleft{}to take a photo\textquotedblright{}.

\begin{grammarlens}[title={What you've built}]
One more word for this chapter.
\end{grammarlens}

\subsection*{Your turn}
\begin{itemize}
\raggedright
  \item \emph{Say it:} \emph{phoṭo kāḍhṇe}
  \item \emph{Say it:} \emph{phoṭo kāḍhṇe}, once more
  \item \emph{Say it:} say \emph{phoṭo kāḍhṇe}
  \item \emph{Recall:} say the Marathi for a mistake, then the Marathi for to repair, then say \emph{phoṭo kāḍhṇe} again
\end{itemize}

\begin{tcolorbox}[breakable,colback=teal!4,colframe=teal!35!black,title={Before you move on}]
What does \mr{फोटो} \mr{काढणे} mean? (\textquotedblleft{}To take a photo\textquotedblright{}.) Say it once more.
\end{tcolorbox}
\section[\texorpdfstring{\mr{समजावणे} (samjāvṇe)}{समजावणे (samjāvṇe)}]{\mr{समजावणे} (samjāvṇe) — to explain}

\label{lesson:MR-C223-samjavne}



\begin{quote}
Before the new one: say the Marathi for to repair, then the Marathi for to take a photo.
\end{quote}

\subsection*{You'll want to know: \mr{समजावणे}}
\textbf{\mr{समजावणे}} — \emph{samjāvṇe} — \textquotedblleft{}to explain\textquotedblright{}.

\begin{grammarlens}[title={What you've built}]
One more word for this chapter.
\end{grammarlens}

\subsection*{Your turn}
\begin{itemize}
\raggedright
  \item \emph{Say it:} \emph{samjāvṇe}
  \item \emph{Say it:} \emph{samjāvṇe}, once more
  \item \emph{Say it:} say \emph{samjāvṇe}
  \item \emph{Recall:} say the Marathi for to repair, then the Marathi for to take a photo, then say \emph{samjāvṇe} again
\end{itemize}

\begin{tcolorbox}[breakable,colback=teal!4,colframe=teal!35!black,title={Before you move on}]
What does \mr{समजावणे} mean? (\textquotedblleft{}To explain\textquotedblright{}.) Say it once more.
\end{tcolorbox}
\section[\texorpdfstring{\mr{भाषांतर} \mr{करणे} (bhāṣāntar karṇe)}{भाषांतर करणे (bhāṣāntar karṇe)}]{\mr{भाषांतर} \mr{करणे} (bhāṣāntar karṇe) — to translate}

\label{lesson:MR-C223-bhasantar-karne}



\begin{quote}
Before the new one: say the Marathi for to take a photo, then the Marathi for to explain.
\end{quote}

\subsection*{You'll want to know: \mr{भाषांतर} \mr{करणे}}
\textbf{\mr{भाषांतर} \mr{करणे}} — \emph{bhāṣāntar karṇe} — \textquotedblleft{}to translate\textquotedblright{}.

\begin{grammarlens}[title={What you've built}]
One more word for this chapter.
\end{grammarlens}

\subsection*{Your turn}
\begin{itemize}
\raggedright
  \item \emph{Say it:} \emph{bhāṣāntar karṇe}
  \item \emph{Say it:} \emph{bhāṣāntar karṇe}, once more
  \item \emph{Say it:} say \emph{bhāṣāntar karṇe}
  \item \emph{Recall:} say the Marathi for to take a photo, then the Marathi for to explain, then say \emph{bhāṣāntar karṇe} again
\end{itemize}

\begin{tcolorbox}[breakable,colback=teal!4,colframe=teal!35!black,title={Before you move on}]
What does \mr{भाषांतर} \mr{करणे} mean? (\textquotedblleft{}To translate\textquotedblright{}.) Say it once more.
\end{tcolorbox}
\section[\texorpdfstring{\mr{निर्णय} \mr{घेणे} (nirṇay gheṇe)}{निर्णय घेणे (nirṇay gheṇe)}]{\mr{निर्णय} \mr{घेणे} (nirṇay gheṇe) — to make a decision}

\label{lesson:MR-C223-nirnay-ghene}



\begin{quote}
Before the new one: say the Marathi for to explain, then the Marathi for to translate.
\end{quote}

\subsection*{You'll want to know: \mr{निर्णय} \mr{घेणे}}
\textbf{\mr{निर्णय} \mr{घेणे}} — \emph{nirṇay gheṇe} — \textquotedblleft{}to make a decision\textquotedblright{}.

\begin{grammarlens}[title={What you've built}]
One more word for this chapter.
\end{grammarlens}

\subsection*{Your turn}
\begin{itemize}
\raggedright
  \item \emph{Say it:} \emph{nirṇay gheṇe}
  \item \emph{Say it:} \emph{nirṇay gheṇe}, once more
  \item \emph{Say it:} say \emph{nirṇay gheṇe}
  \item \emph{Recall:} say the Marathi for to explain, then the Marathi for to translate, then say \emph{nirṇay gheṇe} again
\end{itemize}

\begin{tcolorbox}[breakable,colback=teal!4,colframe=teal!35!black,title={Before you move on}]
What does \mr{निर्णय} \mr{घेणे} mean? (\textquotedblleft{}To make a decision\textquotedblright{}.) Say it once more.
\end{tcolorbox}
